\scanpage{47}
\noindent\hfill atkin \deleted{6} \handwritten{\mspage{7}}

\noindent $=b(11)=1$,$b(77)=2$,and $b(154)=2$; also $b(757)=2$ and $b(5299=7\mathbin{.}757)=3$. We knew already from \handwritten{(3.2)} that if $(m,n)=1$ and $b(m)\handwritten{\geq}1$ then $b(mn)\handwritten{\geq}1$, but the above instances suggested immediately a multiplicative property modulo at least $13^{3}$. In fact a few further hand calculations using the values of $a(n)$ as well as $b(n)$ led us to the following conjecture:
\[
\begin{aligned}
&\text{\uline{Let}}\quad t(n)=c(13^{6}n).\quad \text{\uline{Then if}}\quad (m,n)=1\quad\text{\uline{we have}}\\
&\qquad t(1)t(mn)\handwritten{\equiv}t(m)t(n)\quad(\text{\uline{mod}}\;13^{6})\quad .
\end{aligned}
\tag*{\handwritten{(3.8)}}
\]

\noindent Since $c(13^{6})$ is prime to $13$,we may properly divide by it in the residue class ring modulo $13^{6}$; rewriting $t(n)=c(13^{6}n)/c(13^{6})$ we have \handwritten{(3.8)} with $t(1)=1$,and it was now natural to look for the precise detail analogous to \handwritten{(3.2)}. The natural power of $p$ in the third term is $p^{-1}$,and since $p^{11}\handwritten{\equiv}p^{-1}\;(\mathrm{mod}\;13)$ this agrees with \handwritten{(3,2)}. Thus finally,replacing $6$ by a general $\alpha$,we conjectured the following:
\[
\begin{aligned}
&\qquad\text{\uline{Let}}\quad\alpha\handwritten{\geq}1,\text{\uline{and}}\quad
p\handwritten{\ne}13\quad\text{\uline{be prime}. \uline{Then if}}\quad
t(n)=c(13^{\alpha}n\deleted{\unsure{-|=}}\handwritten{)}/c\handwritten{(}13^{\alpha}\deleted{\unsure{-|=}}\handwritten{)}\\
&\text{\uline{we} \uline{have}, \uline{for} \uline{all}}\quad n\handwritten{\geq}1\handwritten{,}\\
&\qquad t(np)-t(n)t(p)+p^{-1}t(n/p)\handwritten{\equiv}0\quad(\mathrm{mod}\;13^{\alpha})\deleted{\unsure{8|\beta}}\,.
\end{aligned}
\tag*{\handwritten{(3.9)}}
\]
\margin{\handwritten{/\deleted{8}}}

This conjec\rlap{\handwritten{\tikz[baseline=0pt,x=1em,y=1ex]{\draw (0,-.35)--(.56,-.35)--(.56,1.7)--(1.06,1.7);}}}utre I tested successfully on Atlas for $\alpha\handwritten{\leq}8$, $p\handwritten{\leq}97$,and $pn\handwritten{\leq}1000$,and have recently proved in general. We may now regard \handwritten{(3.7)} as the ramified form which \handwritten{(3.9)} takes for the exceptional prime $p=13$,and in fact we can rewrite \handwritten{(3.7)} as
\[
t(13n)\mathrel{\rlap{\raisebox{-.12ex}{\handwritten{\rule{.65em}{.55pt}}}}=}t(13)t(n)\quad(\mathrm{mod}\;13^{\alpha}).
\tag*{\handwritten{(3.10)\tikz[baseline=0pt,x=.18em,y=.6ex]{\fill (0,0) circle (.13);\fill (.65,.8) circle (.13);}}}
\]
We shall \deleted{\unsure{n|m}}in future refer to multiplicative co\deleted{m}\inserted{\handwritten{n}}\margin{\handwritten{/}}gruence properties similar to \handwritten{(3.9)} as M-results,and to their ramified cases similar to \handwritten{(3.10)} as R-results.

Now it is easy to show\rlap{\handwritten{\tikz[baseline=-.2ex,x=.3em,y=.65ex,line cap=round]{\draw (0,0)--(1.2,3.3);\draw (.36,.84)--(.97,.84).. controls (.92,.42) and (1.0,.17)..(1.12,.1);}}}\inserted{\handwritten{from (3.2)}}\margin{\handwritten{\tikz[baseline=-.2ex,x=.3em,y=.65ex,line cap=round]{\draw (0,0)--(1.2,3.3);\draw (.36,.84)--(.97,.84).. controls (.92,.42) and (1.0,.17)..(1.12,.1);}}} by the theory of second order linear difference equations that for each $p\handwritten{\ne}13$ there exists some $\theta$ such that $c(13\mathbin{.}p^{\theta})\handwritten{\equiv}0\;(\mathrm{mod}\;13)$. Thus from \handwritten{(3.9)} we have for \unsure{a}ny $\alpha\handwritten{\geq}1$ the result
\[
\handwritten{c(13^{\alpha}\mathbin{.}P\mathbin{.}n)\equiv0\;(\mathrm{mod}\;13^{\alpha})\quad\text{\uline{if}}\quad(n,P)=1}
\tag*{\handwritten{(3\unsure{,.|..}11)}}
\]

% archive sheet number: 46
